\documentclass[12pt]{letter}

\usepackage[margin=1in]{geometry}
\usepackage{setspace}
\usepackage[colorlinks=true,urlcolor=black]{hyperref}

\signature{Wang Weijia\\School of Economics and Finance\\Xi'an Jiaotong University\\\href{mailto:wangweijia@stu.xjtu.edu.cn}{wangweijia@stu.xjtu.edu.cn}}
\address{Wang Weijia\\School of Economics and Finance\\Xi'an Jiaotong University\\28 Xianning West Road\\Xi'an, Shaanxi 710061\\China\\\href{mailto:wangweijia@stu.xjtu.edu.cn}{wangweijia@stu.xjtu.edu.cn}}
\date{June 9, 2026}

\begin{document}

\begin{letter}{Editors\\\textit{Technology in Society}}

\opening{Dear Editors,}

I am pleased to submit the manuscript ``Government Certification as Technology Governance: Innovation Signals and Institutional Voids in China'' for consideration as a research article in \textit{Technology in Society}. Following the transfer recommendation from \textit{Technological Forecasting and Social Change}, I have revised the manuscript substantially to fit the aims and readership of \textit{Technology in Society}.

The revised paper studies how a public technology certification changes the social recognition of innovation capability. I focus on China's ``Specialized, Refined, Distinctive, and Innovative'' (SRDI) certification program, a national screening system for technology-oriented small and medium-sized enterprises. The paper argues that government certification should be understood not only as industrial-policy support, but also as public signal infrastructure: it helps banks, customers, suppliers, and local governments identify firms whose technological capability is difficult to observe.

The manuscript fits \textit{Technology in Society} because it links technological upgrading to social and institutional context. The central question is whether a national technology certification works differently across local institutional environments. Using A-share firm data from 2007 to 2021, hand-collected certification records, patent-quality measures, and staggered difference-in-differences designs, the paper shows that certification relaxes financing constraints, increases targeted investment, and promotes labor specialization. The evidence on final innovation quality is more cautious: heterogeneity-robust estimates suggest larger gains in lower-institution cities than in higher-institution cities, but those subgroup estimates are statistically imprecise. I therefore frame the findings as suggestive quasi-experimental evidence rather than as definitive causal proof.

I have made six substantive revisions for this submission. First, I reframed the paper around technology governance, public certification, and institutional voids, with the addition of dedicated citations to the institutional-voids literature (North 1990; Khanna and Palepu 1997; Mair and Marti 2009). Second, I strengthened the theoretical framework by incorporating an information-asymmetry signaling foundation (Akerlof 1970; Leland and Pyle 1977; Stiglitz and Weiss 1981), adding a two-type signaling logic that generates the institutional-substitution prediction, and connecting the broader industrial-policy evaluation literature (Rodrik 2004; Aghion et al.\ 2015; Lerner 2009). Third, I revised the abstract, introduction, discussion, and conclusion to foreground the social mechanisms through which certification changes technology recognition and innovation pathways. Fourth, I softened causal language where the evidence is statistically fragile, especially around pre-treatment dynamics, province-level clustering, and the weak-instrument diagnostic, and added a dedicated subsection explaining the methodological source of divergence between TWFE and heterogeneity-robust estimates. Fifth, I reorganized the submission into clean separate files: an anonymous manuscript, a title page with author details, declarations, author biography, and this cover letter. Sixth, I added province-level wild-cluster bootstrap inference (Webb weights, 9,999 replications) for the preferred interaction specification, reporting bootstrap $p$-values in the main text and table footnote, and redesigned the placebo permutation figure as a histogram and kernel density overlay with the observed estimate marked by a dashed line, making the permutation evidence more transparent.

The manuscript is not under review at any other journal. I declare no financial or personal conflicts of interest. The research received no external grant funding. The study uses firm-level archival data and public policy records and does not involve human subjects. Replication code and processed nonrestricted materials can be made available upon request during review and upon acceptance; access to licensed CSMAR data is subject to the database provider's terms.

Thank you for considering this manuscript. I would be grateful for the opportunity to receive the journal's assessment of the revised paper.

\closing{Yours sincerely,}

\end{letter}
\end{document}
